\documentclass[11pt,a4paper]{article}

\usepackage{amsmath,amssymb,amsthm}
\usepackage{booktabs}
\usepackage{array}
\usepackage{geometry}
\usepackage{parskip}
\usepackage{xcolor}
\usepackage{hyperref}
\usepackage{cleveref}
\usepackage{microtype}
\usepackage{enumitem}

\geometry{margin=2.5cm}

\theoremstyle{definition}
\newtheorem{definition}{Definition}[section]
\newtheorem{example}[definition]{Example}

\theoremstyle{plain}
\newtheorem{theorem}[definition]{Theorem}
\newtheorem{proposition}[definition]{Proposition}
\newtheorem{lemma}[definition]{Lemma}

\theoremstyle{remark}
\newtheorem{remark}[definition]{Remark}

% Notation shorthands
\newcommand{\R}{\mathbb{R}}
\newcommand{\C}{\mathbb{C}}
\newcommand{\N}{\mathbb{N}}
\newcommand{\Sym}{\mathfrak{S}}
\newcommand{\eval}{\mathrm{eval}}
\newcommand{\fo}{\mathrm{FO}}
\newcommand{\repL}{\mathrm{repL}}
\newcommand{\sign}{\mathrm{sign}}
\newcommand{\pf}{\mathrm{pf}}
\newcommand{\acount}{\mathrm{count}}
\newcommand{\comb}[2]{\binom{#1}{#2}}

\title{\textbf{The Symcoder Encoder: Methodology and Correctness}\\
\large A pedagogical account of the permutation-invariant embedding algorithm}
\author{}
\date{}

\begin{document}

\maketitle
\tableofcontents
\newpage

% ---------------------------------------------------------------------------
\section{Overview}
% ---------------------------------------------------------------------------

The encoder takes as input a \emph{physics event}: a finite collection of
particle vectors grouped by particle type, together with a set of scalar
operations.  It produces a fixed-length real vector that is invariant under any
permutation of particles of the same type, and that contains enough information
to separate generic events.

The construction proceeds in two phases.
\begin{enumerate}
  \item \textbf{Phase 1 — single-orbit encoding.}  For each distinct
    \emph{FlavouredOperator} (an operation with a fixed recipe for drawing
    argument labels from the particle groups), evaluate the operation on every
    valid label combination and sort the results.  This yields a
    permutation-invariant sequence of real numbers describing the marginal
    distribution of each scalar feature.

  \item \textbf{Phase 2 — pair-association encoding.}  For each pair of
    FlavouredOperators, evaluate the complex number $z_k = \eval(u_k) +
    i\,\eval(v_k)$ for every pair $(u_k,v_k)$ of atoms that share a specific
    pattern of label overlap (\emph{one association}).  Embed the resulting
    multiset into permutation-invariant polynomial coefficients.  Several
    optimisations reduce storage without losing information.
\end{enumerate}

The two phases together form a complete, permutation-invariant encoding.
Phase~1 captures individual feature statistics; Phase~2 captures pairwise
correlations between features.

% ---------------------------------------------------------------------------
\section{Mathematical Setup}
% ---------------------------------------------------------------------------

\subsection{Events, Groups, and Labels}

\begin{definition}[Vector group]
A \emph{vector group} $\mathcal{V} = (\mathcal{L}, d)$ consists of a finite
label set $\mathcal{L} = \{\ell_1,\ldots,\ell_n\}$ and a dimension $d \in
\N$.  The \emph{size} of the group is $|\mathcal{L}| = n$.
\end{definition}

\begin{definition}[Context]
A \emph{context} $\mathcal{C} = (\mathcal{V}_1,\ldots,\mathcal{V}_m)$ is an
ordered tuple of vector groups with disjoint label sets.  Let $n_i =
|\mathcal{V}_i|$ be the size of group~$i$.
\end{definition}

\begin{definition}[Event]
An \emph{event} $E$ is an assignment $\ell \mapsto \mathbf{v}_\ell \in \R^d$
of a $d$-dimensional vector to every label $\ell$ across all groups.
\end{definition}

The symmetry group acting on events is the product of symmetric groups
\[
  G = \Sym_{n_1} \times \cdots \times \Sym_{n_m},
\]
where each $\Sym_{n_i}$ permutes the labels of group~$i$.  Two events are
\emph{physically equivalent} if one is obtained from the other by a $G$-action,
i.e.\ by relabelling particles within each type.  The encoder must produce the
same output for physically equivalent events.

\subsection{Operations}

\begin{definition}[Operation]
An \emph{operation} $\mathit{op}$ is characterised by:
\begin{itemize}[nosep]
  \item a \emph{rank} $r \in \N$ (number of vector arguments),
  \item a \emph{parity} $p \in \{+1,-1\}$,
  \item an \emph{argument symmetry} $\sigma \in \{\textsc{Symmetric},
    \textsc{Antisymmetric}\}$, and
  \item an \emph{evaluation function} $\mathit{op}_{fn}: (\R^d)^r \to \R$.
\end{itemize}
For \textsc{Symmetric} operations, $\mathit{op}_{fn}$ is invariant under all
permutations of its $r$ arguments.  For \textsc{Antisymmetric} operations,
$\mathit{op}_{fn}$ changes sign under any transposition of two arguments.
\end{definition}

\begin{example}
The dot product $\mathit{dot}(\mathbf{u},\mathbf{v}) = \mathbf{u}\cdot\mathbf{v}$ has
rank~2, parity~$+1$, and is \textsc{Symmetric}.

The triple product
$\varepsilon_3(\mathbf{u},\mathbf{v},\mathbf{w}) =
\mathbf{u}\cdot(\mathbf{v}\times\mathbf{w})$ has rank~3, parity~$-1$, and is
\textsc{Antisymmetric}.
\end{example}

\subsection{Atoms}

An \emph{atom} is an operation applied to a specific ordered tuple of labels.
Because the argument symmetry of the operation constrains how permutations of
the labels affect the result, each unordered combination of labels gives rise to
either one or two atoms.

\begin{definition}[Atom]
Let $\mathit{op}$ be an operation of rank~$r$ and let
$(\ell_1,\ldots,\ell_r)$ be a tuple of distinct labels in canonical
(lexicographic) order within each group.  An \emph{atom} is a triple
\[
  u = (\mathit{op},\, (\ell_1,\ldots,\ell_r),\, \varepsilon)
\]
where $\varepsilon \in \{+1,-1\}$ is the \emph{sign}.

\begin{itemize}
  \item For \textsc{Symmetric} $\mathit{op}$: only $\varepsilon = +1$ arises.
  \item For \textsc{Antisymmetric} $\mathit{op}$: both $\varepsilon = +1$ and
    $\varepsilon = -1$ arise, one for each canonical label tuple.  The sign
    $\varepsilon$ records the parity of the permutation that maps the canonical
    ordering to a particular chosen reference ordering.
\end{itemize}
\end{definition}

\begin{definition}[Evaluation]
The evaluation of atom $u = (\mathit{op}, (\ell_1,\ldots,\ell_r), \varepsilon)$
on event $E$ is
\[
  \eval(u, E) = \varepsilon \cdot \mathit{op}_{fn}(\mathbf{v}_{\ell_1},
  \ldots, \mathbf{v}_{\ell_r}).
\]
\end{definition}

The sign $\varepsilon$ means that for an \textsc{Antisymmetric} operation the
two atoms sharing the same label tuple satisfy
$\eval(u^{+}, E) = -\eval(u^{-}, E)$ for every event $E$, by antisymmetry of
$\mathit{op}_{fn}$.

\subsection{Flavours and FlavouredOperators}

Atoms are organised by how many labels each argument draws from each group.

\begin{definition}[Flavour]
A \emph{flavour} for a context with $m$ groups is a tuple $F = (k_1,\ldots,k_m)$
of non-negative integers with $\sum_i k_i = r$ (the rank of the associated
operation).  We require $k_i \le n_i$ for all~$i$.
\end{definition}

\begin{definition}[FlavouredOperator]
A \emph{FlavouredOperator} $\fo(\mathit{op}, F)$ is the set of all atoms whose
operation is $\mathit{op}$ and whose label tuple draws exactly $k_i$ labels from
group~$i$, for each~$i$.
\end{definition}

The number of atoms in a FlavouredOperator is
\begin{equation}\label{eq:fo-count}
  |\fo(\mathit{op}, F)| =
  \begin{cases}
    \displaystyle\prod_{i=1}^m \comb{n_i}{k_i}
      & \text{if $\mathit{op}$ is \textsc{Symmetric}},\\[8pt]
    \displaystyle 2\prod_{i=1}^m \comb{n_i}{k_i}
      & \text{if $\mathit{op}$ is \textsc{Antisymmetric}.}
  \end{cases}
\end{equation}
The factor of~2 for \textsc{Antisymmetric} operations reflects the two sign
choices per label combination.

\begin{definition}[repL]
For a context $\mathcal{C}$ and a set of operations $\mathcal{K}$, the
\emph{repL} is the list of all valid FlavouredOperators:
\[
  \repL(\mathcal{C}, \mathcal{K})
  = \bigl\{\fo(\mathit{op}, F) : \mathit{op} \in \mathcal{K},\;
    F \text{ is a valid flavour for $\mathit{op}$ and $\mathcal{C}$}\bigr\}.
\]
\end{definition}

Every atom in the encoding vocabulary appears in exactly one FlavouredOperator.

% ---------------------------------------------------------------------------
\section{Phase 1: Single-Orbit Encoding}
% ---------------------------------------------------------------------------

The first phase encodes each FlavouredOperator's evaluation values as a
permutation-invariant sequence of reals.

\subsection{The Symmetric Case}

For a \textsc{Symmetric} FlavouredOperator $\mathit{fo} = \fo(\mathit{op}, F)$,
let $n = \prod_i \comb{n_i}{k_i}$ be the number of atoms (all with sign
$+1$).  Evaluating on event $E$ gives $n$ real numbers
\[
  a_1 = \eval(u_1, E),\quad a_2 = \eval(u_2, E),\quad \ldots,\quad
  a_n = \eval(u_n, E).
\]
The \textbf{Phase~1 encoding} of $\mathit{fo}$ is the sorted sequence
\[
  \bigl(a_{(1)},\, a_{(2)},\, \ldots,\, a_{(n)}\bigr),
  \qquad a_{(1)} \le a_{(2)} \le \cdots \le a_{(n)}.
\]
Sorting is the canonical permutation-invariant operation.  The sorted sequence
is a complete encoding of the multiset $\{a_1,\ldots,a_n\}$, and hence of all
symmetric functions of that multiset.

\subsection{The Antisymmetric Case: Sign Compression}
\label{sec:5c}

For an \textsc{Antisymmetric} $\mathit{fo}$, the atom set contains pairs
$(u^+, u^-)$ for each label combination, with $\eval(u^+, E) =
-\eval(u^-, E)$.  The full set of $2n$ evaluation values is therefore
\[
  \{a_1, -a_1, a_2, -a_2, \ldots, a_n, -a_n\},
  \qquad \text{where } a_j = \eval(u^+_j, E).
\]
This multiset is antisymmetric about zero: knowing the positive half
$\{a_1,\ldots,a_n\}$ (or equivalently the absolute values $\{|a_j|\}$)
determines the full multiset.

\begin{proposition}[Sign compression]
The multiset $\{|a_1|,\ldots,|a_n|\}$ is invariant under the action of~$G$.
\end{proposition}
\begin{proof}
A group element $g \in G$ permutes atoms and may change the canonical
label ordering, thereby flipping the sign of some atoms (i.e.\ swapping $u^+_j
\leftrightarrow u^-_j$ for certain $j$).  However, $|a_j| = |\eval(u^+_j,E)|
= |\eval(u^-_j,E)|$ is invariant under sign flip, so the multiset
$\{|a_j|\}$ is unchanged.
\end{proof}

The decision to apply sign compression is made from the operation's argument
symmetry alone — never from the numerical values of $a_j$.

\textbf{Phase~1 encoding (antisymmetric case):} store the sorted sequence
\[
  \bigl(|a|_{(1)}, |a|_{(2)}, \ldots, |a|_{(n)}\bigr),
  \qquad |a|_{(1)} \le \cdots \le |a|_{(n)},
\]
contributing $n = \frac{1}{2}|\mathit{fo}|$ reals instead of $2n$.

\subsection{Summary}

Each distinct FlavouredOperator $\mathit{fo}$ contributes
\begin{equation}\label{eq:phase1-length}
  \ell_1(\mathit{fo}) =
  \begin{cases}
    |\mathit{fo}| & \text{if \textsc{Symmetric}},\\
    \tfrac{1}{2}|\mathit{fo}| & \text{if \textsc{Antisymmetric}},
  \end{cases}
\end{equation}
reals to Phase~1.  In both cases $\ell_1(\mathit{fo}) = \prod_i \comb{n_i}{k_i}$,
the number of distinct label combinations.

% ---------------------------------------------------------------------------
\section{Phase 2: Pair-Association Encoding}
% ---------------------------------------------------------------------------

Phase~2 encodes pairwise correlations between atoms from (possibly different)
FlavouredOperators.  The central objects are \emph{PairFlavours}, which classify
atom-pairs by orbit type under $G$.

\subsection{PairFlavours and Overlap Blocks}

\begin{definition}[PairFlavour]
A \emph{PairFlavour} is a tuple
\[
  \pf = (\mathit{op}_u, F_u, \mathit{op}_v, F_v, s),
\]
where $(op_u, F_u)$ and $(\mathit{op}_v, F_v)$ are the operation-flavour pairs
of atoms $u$ and~$v$, and $s = (s_1,\ldots,s_m)$ is the \emph{overlap} vector:
$s_i$ is the number of labels shared between $u$ and $v$ from group~$i$.
We require $0 \le s_i \le \min(k_{u,i}, k_{v,i})$ for all~$i$.

By convention the two sides are canonically ordered so that $(\mathit{op}_u,
F_u) \le (\mathit{op}_v, F_v)$ under a fixed total order; this ensures that
$\pf(u,v) = \pf(v,u)$.
\end{definition}

\begin{proposition}
Two atom-pairs $(u,v)$ and $(u',v')$ lie in the same $G$-orbit if and only if
they have the same PairFlavour.
\end{proposition}

\begin{definition}[Overlap block]
An \emph{overlap block} is the set of all PairFlavours sharing the same
$(\mathit{op}_u, F_u, \mathit{op}_v, F_v)$, differing only in the overlap~$s$.
Each member of an overlap block is called an \emph{association}.

The valid overlap range for group~$i$ is
\[
  s_i \in \bigl[\max(0,\, k_{u,i}+k_{v,i}-n_i),\; \min(k_{u,i}, k_{v,i})\bigr].
\]
\end{definition}

\subsection{The Association Table}

Fix an overlap block $B = \{(\mathit{op}_u, F_u, \mathit{op}_v, F_v, s) : s
\text{ valid}\}$.  Imagine a matrix whose columns are labelled by the atoms of
$\fo(\mathit{op}_u, F_u)$ and whose rows are labelled by the atoms of
$\fo(\mathit{op}_v, F_v)$.  Each cell $(u_j, v_k)$ belongs to exactly one
association: the one whose overlap $s_i$ equals the number of labels shared by
$u_j$ and $v_k$ from group~$i$.

\begin{example}
\label{ex:assoc-table}
Let $n=3$ (labels $a,b,c$), one group, and $\mathit{op} = \mathit{dot}$ with
rank~2 and flavour $F=(2)$.  The atoms are $\mathit{dot}(a,b)$,
$\mathit{dot}(a,c)$, $\mathit{dot}(b,c)$.  The self-pairing overlap block
$(\mathit{dot},(2),\mathit{dot},(2))$ has associations $s \in \{1,2\}$ (overlap
$s=0$ is impossible here since $k_u+k_v-n = 2+2-3 = 1$).

\medskip
\begin{center}
\renewcommand{\arraystretch}{1.4}
\begin{tabular}{c|ccc}
  \toprule
  & $\mathit{dot}(a,b)$ & $\mathit{dot}(a,c)$ & $\mathit{dot}(b,c)$ \\
  \midrule
  $\mathit{dot}(a,b)$ & $s=2$ & $s=1$ & $s=1$ \\
  $\mathit{dot}(a,c)$ & $s=1$ & $s=2$ & $s=1$ \\
  $\mathit{dot}(b,c)$ & $s=1$ & $s=1$ & $s=2$ \\
  \bottomrule
\end{tabular}
\end{center}
\medskip

Association $s=2$ (diagonal): 3 cells; association $s=1$ (off-diagonal): 6 cells.
\end{example}

The number of atom-pairs in association $\pf$ is
\begin{equation}\label{eq:pf-count}
  \acount(\pf) = \prod_{i=1}^m \comb{n_i}{s_i}\comb{n_i-s_i}{k_{u,i}-s_i}
    \comb{n_i-k_{u,i}}{k_{v,i}-s_i}.
\end{equation}
The three binomial factors count: (i)~the shared labels, (ii)~the extra
$u$-labels from the remaining pool, (iii)~the extra $v$-labels from what is
left.

\begin{remark}
The cells of the association table partition exactly into the associations of
the block:
\[
  |\fo(\mathit{op}_u,F_u)| \cdot |\fo(\mathit{op}_v,F_v)|
  = \sum_{\pf \in B} \acount(\pf)
\]
(counting sign variants on both sides for antisymmetric operations).
\end{remark}

\subsection{Dropping the Largest Association: The Complementation Argument}
\label{sec:5b}

Within each overlap block, the association with the largest $\acount(\pf)$ is
not encoded explicitly.  Instead, it is recovered by the decoder via a
multiset-complementation argument.

Let $B$ be an overlap block, and let $\pf^* \in B$ be the association with the
largest count.  Write
\[
  U = \{\eval(u_j, E) : u_j \in \fo(\mathit{op}_u, F_u)\}
  \quad\text{and}\quad
  V = \{\eval(v_k, E) : v_k \in \fo(\mathit{op}_v, F_v)\}
\]
for the multisets of evaluation values (known from Phase~1 encodings of
$\fo(\mathit{op}_u, F_u)$ and $\fo(\mathit{op}_v, F_v)$).

The multiset of complex numbers
\[
  Z = \{z_{jk} = \eval(u_j,E) + i\,\eval(v_k,E) : (u_j,v_k)
  \text{ any pair with $u_j \in \fo_u$, $v_k \in \fo_v$}\}
\]
is the multiset product $U \oplus_\C V = \{a + ib : a \in U, b \in V\}$, which
is determined entirely by $U$ and~$V$.  The associations partition $Z$ into
disjoint submultisets $Z_\pf$ for $\pf \in B$.  Therefore, knowing all
associations except $\pf^*$ determines $Z_{\pf^*}$ by complement:
\[
  Z_{\pf^*} = Z \;\setminus\; \bigsqcup_{\pf \in B,\, \pf\ne \pf^*} Z_\pf.
\]

\begin{proposition}[Complementation is exact]
The above complementation works for \emph{all} events $E$, including those
where some evaluation values coincide.  Since $Z$ is a multiset, repeated
values appear with the correct multiplicity, and the complement operation on
multisets accounts for this correctly.
\end{proposition}

\begin{remark}[Choosing which association to drop]
The complementation argument is valid regardless of which association within the
block is designated as dropped.  Choosing the largest minimises storage, since
the number of reals contributed by an association scales with $\acount(\pf)$.
\end{remark}

\begin{remark}[Single-association blocks]
If an overlap block contains only one association (which occurs when
$\fo(\mathit{op}_u, F_u)$ and $\fo(\mathit{op}_v, F_v)$ draw labels from
completely disjoint groups, forcing $s = (0,\ldots,0)$ as the only valid
overlap), then the entire block is dropped: it is the unique and hence largest
association in its block.  No information is lost, since the single
association's $z$-values are a Cartesian product of $U$ and $V$, which are
already encoded in Phase~1.
\end{remark}

\subsection{The Z-value Embedding}
\label{sec:z-embed}

For each surviving association $\pf$ (i.e.\ each non-dropped member of each
overlap block), the encoder computes a set of complex numbers and then produces
a permutation-invariant polynomial encoding.

\subsubsection{The positive-sign subset}

For each atom-pair $(u_j, v_k)$ with both $u_j.\varepsilon = +1$ and
$v_k.\varepsilon = +1$, define
\[
  z_k = \eval(u_j, E) + i\,\eval(v_k, E).
\]
This subset has exactly $\acount(\pf)$ elements regardless of which operations
are antisymmetric, because for each base label combination there is exactly one
$(u^+, v^+)$ pair.

\subsubsection{Permutation-invariant embedding}

The multiset $\{z_k\}$ (of size $n = \acount(\pf)$) is embedded by the
Echinocoder polynomial encoder: compute the monic polynomial whose roots are
$\{-z_k\}$, and return all coefficients except the leading (monic) one,
ordered by descending degree.  This yields $n$ complex numbers that are
invariant under any permutation of the $z_k$ values (i.e.\ under any relabelling
of the pairs).

Each complex coefficient is unpacked into two reals $(\mathrm{Re}, \mathrm{Im})$,
contributing $2n = 2\acount(\pf)$ reals in total.

\subsection{Sign-Symmetry Compression of the Embedding}
\label{sec:5a}

When one or both operations in an association are \textsc{Antisymmetric}, the
full orbit of $z$-values (over all sign combinations of $u$ and $v$) has
a forced structure that can be exploited to halve (or quarter) the apparent cost
of embedding the full orbit.

The key observation is that the full orbit's $z$-values come in symmetry-related
groups:
\begin{itemize}
  \item $\textsc{Sym}\times\textsc{Antisym}$ (\emph{SA}):
    the full orbit is $\{z_k, \overline{z}_k\}$ — conjugate-closed.
  \item $\textsc{Antisym}\times\textsc{Sym}$ (\emph{AS}):
    the full orbit is $\{z_k, -\overline{z}_k\}$ — anti-conjugate-closed.
  \item $\textsc{Antisym}\times\textsc{Antisym}$ (\emph{AA}):
    the full orbit contains $\{z_k, \overline{z}_k, -z_k, -\overline{z}_k\}$,
    so the squares $\{z_k^2\}$ form a conjugate-closed set of size~$n$.
\end{itemize}

In each non-trivial case, the characteristic polynomial of the full orbit
has forced coefficient structure (real coefficients, or alternating
real/imaginary, etc.) that allows $2n$ reals to encode the same information as
the brute-force $4n$ or $8n$ reals.  The details are as follows.

\subsubsection{SS: Symmetric $\times$ Symmetric}
Embed $\{z_k\}$ ($n$ complex roots) directly.  Output: $n$ complex
$= 2n$ reals.  No forced structure.

\subsubsection{SA: Symmetric $\times$ Antisymmetric}
The full orbit is $\{z_k, \overline{z}_k\}$ ($2n$ roots).  The corresponding
polynomial has \emph{real} coefficients (a conjugate-closed root set implies
real coefficients).  Output: $2n$ real coefficients, packed as $n$ complex
numbers ($2n$ reals total).

\subsubsection{AS: Antisymmetric $\times$ Symmetric}
The full orbit is $\{z_k, -\overline{z}_k\}$ ($2n$ roots).  The polynomial has
even-degree coefficients real and odd-degree coefficients purely imaginary.
Output: $n$ independent reals of each parity, again packed as $n$ complex
numbers ($2n$ reals).

\subsubsection{AA: Antisymmetric $\times$ Antisymmetric}
The full orbit is $\{z_k, \overline{z}_k, -z_k, -\overline{z}_k\}$ ($4n$
roots in $n$ quadruples).  Setting $w_k = z_k^2$ reduces to an SA-type
structure in the variable~$w$: the multiset $\{w_k, \overline{w}_k\}$ is
conjugate-closed and its polynomial has real coefficients.  Output: $n$ complex
$= 2n$ reals, after embedding $\{w_k^2\}$ as in SA.

\begin{proposition}[Output size]\label{prop:assoc-size}
Every association, regardless of symmetry class, contributes exactly
$2\acount(\pf)$ reals to the Phase~2 encoding.
\end{proposition}

The compression is in avoiding the naive approach of embedding the full orbit
(which would require $2 \times \mathrm{orbit\_size}(\pf)$ reals and grows by a
factor of~2 or~4 for antisymmetric operations).

% ---------------------------------------------------------------------------
\section{Correctness}
% ---------------------------------------------------------------------------

\subsection{Permutation Invariance}

\begin{theorem}
The complete encoding (Phase~1 + Phase~2) is invariant under the action of the
symmetry group $G = \Sym_{n_1} \times \cdots \times \Sym_{n_m}$.
\end{theorem}

\begin{proof}[Proof sketch]
\textbf{Phase~1.}  A group element $g \in G$ permutes atoms within each
FlavouredOperator, and for \textsc{Antisymmetric} operations may flip the sign
of some atoms.  The \textsc{Symmetric} Phase~1 encoding sorts evaluation values,
and sorting commutes with any permutation of the values.  The
\textsc{Antisymmetric} encoding sorts absolute values, which are additionally
invariant under sign flips (\cref{sec:5c}).

\textbf{Phase~2.}  The complex numbers $z_k = \eval(u_j,E) +
i\,\eval(v_k,E)$ are defined for pairs $(u_j,v_k)$ in the positive-sign subset
of association $\pf$.  Under $g$, the set of such pairs is permuted (with
possible sign changes, handled by the 5a symmetry classes in
\cref{sec:5a}).  The polynomial embedding of a multiset is invariant under any
permutation of its elements, so the embedding is $G$-invariant.  The
complementation step (\cref{sec:5b}) produces the same dropped-association
values because it operates on the multiset $Z_{\pf^*}$, which is the
set-theoretic complement of the stored associations within $Z = U \oplus_\C V$,
and $Z$, $U$, $V$ are all $G$-invariant.
\end{proof}

\subsection{Information Completeness}

The encoding is \emph{complete} if the full set of evaluation values for all
atoms and atom-pairs can be reconstructed from the encoded vector.  Completeness
is required for the encoder to separate generic events.

\begin{theorem}[Phase~1 completeness]
The Phase~1 encoding of $\fo(\mathit{op}, F)$ is a complete encoding of the
multiset $\{\eval(u, E) : u \in \fo\}$.
\end{theorem}
\begin{proof}
The sorted sequence is the canonical representative of the multiset; it uniquely
determines the multiset.  For the \textsc{Antisymmetric} case, the sorted
absolute values $\{|a_j|\}$ determine the multiset $\{a_j, -a_j\}$ because each
absolute value contributes exactly one positive/negative pair.
\end{proof}

\begin{theorem}[Phase~2 completeness]
Given the Phase~1 encodings of all FlavouredOperators, the Phase~2 encoding
(after complementation) is complete.
\end{theorem}
\begin{proof}
For each overlap block $B$:
\begin{enumerate}
  \item The stored associations' $z$-multisets $\{Z_\pf\}_{\pf \ne \pf^*}$
    are recovered from their polynomial embeddings (which are injective on
    multisets by the Echinocoder construction).
  \item The Phase~1 encodings supply the multisets $U$ and $V$, hence $Z = U
    \oplus_\C V$.
  \item The dropped association satisfies $Z_{\pf^*} = Z \setminus
    \bigsqcup_{\pf \ne \pf^*} Z_\pf$, uniquely determined as a multiset
    complement.
  \item The 5a sign-symmetry compression is reversible: the full orbit
    $z$-multiset is reconstructed from the positive-sign subset and the
    symmetry class of the association.
\end{enumerate}
\end{proof}

\begin{remark}[Degenerate events]
Completeness holds for \emph{all} events, including degenerate ones where some
particles have identical vectors.  Identical particles produce repeated
evaluation values; these are handled correctly as multiset multiplicities
throughout (in Phase~1 sorted sequences, in $Z = U \oplus_\C V$, and in the
multiset complement).  Two numerically identical particles are indistinguishable,
which is the physically correct behaviour for a permutation-invariant encoder.
\end{remark}

% ---------------------------------------------------------------------------
\section{Output Structure}
% ---------------------------------------------------------------------------

The output of the encoder for a given plan $(\mathcal{C}, \mathcal{K})$ and
event $E$ is a concatenation of real-valued segments.  The structure depends
only on the plan; it is the same for every event.

\subsection{Segment types}

\begin{description}
  \item[\texttt{ORBIT}] One per distinct FlavouredOperator $\mathit{fo}$ in
    $\repL(\mathcal{C},\mathcal{K})$.  Length: $\prod_i \comb{n_i}{k_i}$ reals
    (\emph{sign-compressed} for \textsc{Antisymmetric} operations; label
    \texttt{sign\_compressed=True} in the metadata).

  \item[\texttt{ASSOC}] One per stored association (every non-dropped member of
    each overlap block).  Length: $2\acount(\pf)$ reals.

  \item[\texttt{NULL}] One per dropped association (largest per block, by
    $\acount(\pf)$).  Length: 0 reals.  Present in the metadata at its logical
    position in the overlap block so a decoder knows which association was
    omitted.
\end{description}

All \texttt{ORBIT} segments precede all \texttt{ASSOC}/\texttt{NULL} segments.
Within Phase~2, segments are grouped by overlap block, ordered by a canonical
sort key $(\mathit{op}_u\text{-key}, \mathit{op}_v\text{-key}, s)$ (where
$\mathit{op}$-key $= (\mathrm{name}, r, p, \sigma, F)$), with overlap~$s$
ascending lexicographically within each block.

\subsection{Total output length}

\begin{equation}
  L = \underbrace{\sum_{\mathit{fo} \in \repL} \prod_i \comb{n_i}{k_i}}_{\text{Phase~1}}
    + \underbrace{\sum_{\substack{\pf \in \text{all blocks} \\ \pf \ne \pf^*_B}}\!
      2\,\acount(\pf)}_{\text{Phase~2}},
\end{equation}
where $\pf^*_B$ denotes the dropped (largest-count) association in block~$B$.

\subsection{Human-readable metadata}

The function \texttt{describe\_encoding(plan)} returns a list of
\texttt{SegmentInfo} records, one per segment, containing: segment kind
(\texttt{ORBIT}/\texttt{ASSOC}/\texttt{NULL}), start index, length, operation
names, flavours, overlap, symmetry class (\texttt{SS}/\texttt{SA}/\texttt{AS}/\texttt{AA}),
and the \texttt{sign\_compressed} flag for \texttt{ORBIT} segments.  This list
is a complete description of the encoder's output structure given only the
plan — no event data is needed.

% ---------------------------------------------------------------------------
\section{Algorithm Summary}
% ---------------------------------------------------------------------------

\begin{center}
\renewcommand{\arraystretch}{1.5}
\begin{tabular}{lll}
  \toprule
  \textbf{Step} & \textbf{Input} & \textbf{Output} \\
  \midrule
  Construct \texttt{repL} & plan $(\mathcal{C}, \mathcal{K})$
    & list of FlavouredOperators \\
  Compute PairFlavours & \texttt{repL}
    & all valid $\pf$, grouped into overlap blocks \\
  \midrule
  Phase~1 per $\mathit{fo}$
    & \texttt{fo.atoms()}, event $E$
    & $\ell_1(\mathit{fo})$ sorted reals \\
  \midrule
  Phase~2: choose $\pf^*$ per block & overlap blocks & dropped-association map \\
  Phase~2: $z$-values per $\pf \ne \pf^*$
    & atom pairs, event $E$
    & $\acount(\pf)$ complex numbers \\
  Phase~2: embed $z$-values (5a)
    & $\{z_k\}$, symmetry class
    & $\acount(\pf)$ complex coefficients \\
  Phase~2: unpack to reals
    & complex coefficients
    & $2\acount(\pf)$ reals \\
  \midrule
  Concatenate & all segments & output vector of length $L$ \\
  \bottomrule
\end{tabular}
\end{center}

% ---------------------------------------------------------------------------
\appendix
\section{Notation Reference}
% ---------------------------------------------------------------------------

\begin{center}
\renewcommand{\arraystretch}{1.4}
\begin{tabular}{ll}
  \toprule
  \textbf{Symbol} & \textbf{Meaning} \\
  \midrule
  $m$ & number of particle groups in the context \\
  $n_i$ & size of group $i$ (number of labels) \\
  $G = \Sym_{n_1}\times\cdots\times\Sym_{n_m}$ & symmetry group \\
  $r$ & rank of an operation (number of vector arguments) \\
  $F = (k_1,\ldots,k_m)$ & flavour: label counts per group, $\sum k_i = r$ \\
  $\fo(\mathit{op}, F)$ & FlavouredOperator \\
  $\varepsilon \in \{+1,-1\}$ & sign of an atom \\
  $\eval(u, E)$ & scalar evaluation of atom $u$ on event $E$ \\
  $\acount(\pf)$ & number of atom-pairs in association $\pf$ \\
  $s = (s_1,\ldots,s_m)$ & overlap vector of a PairFlavour \\
  $Z_\pf$ & multiset of $z$-values for association $\pf$ \\
  $U \oplus_\C V$ & complex Cartesian sum $\{a+ib : a\in U, b\in V\}$ \\
  \bottomrule
\end{tabular}
\end{center}

\end{document}
